% ================================================================
%  ХАВТАСНЫ ХУУДАС (ГАРЧГИЙН ХУУДАС 1)
% ================================================================
\begin{titlepage}
\pagestyle{empty}
\centering

\vspace*{0.5cm}
{\large\bfseries Шинжлэх Ухаан Технологийн Их Сургууль}\\[6pt]
{\large Мэдээлэл, Холбооны Технологийн Сургууль}

\vspace{3cm}
{\Large\bfseries Тогтохын Оюунбаатар}

\vspace{2cm}
{\LARGE\bfseries QR КОДООР СУУРИЛСАН РЕСТОРАНЫ}\\[8pt]
{\LARGE\bfseries ШИРЭЭНИЙ МОБАЙЛ ЗАХИАЛГЫН}\\[8pt]
{\LARGE\bfseries СИСТЕМ ХӨГЖҮҮЛЭХ НЬ}

\vspace{2cm}
{\large\itshape Бакалаврын төгсөлтийн ажил}

\vfill
{\large Улаанбаатар хот}\\[4pt]
{\large 2025 он}
\end{titlepage}
\clearpage

% ================================================================
%  ХАВТАСНЫ ХУУДАС (ГАРЧГИЙН ХУУДАС 2)
% ================================================================
\begin{titlepage}
\pagestyle{empty}
\centering

{\large\bfseries Шинжлэх Ухаан Технологийн Их Сургууль}\\[6pt]
{\large Мэдээлэл, Холбооны Технологийн Сургууль}

\vspace{0.5cm}
\begin{flushright}
{\normalsize Мэдээллийн технологийн салбар}
\end{flushright}

\vspace{1.5cm}
{\LARGE\bfseries QR КОДООР СУУРИЛСАН РЕСТОРАНЫ}\\[8pt]
{\LARGE\bfseries ШИРЭЭНИЙ МОБАЙЛ ЗАХИАЛГЫН}\\[8pt]
{\LARGE\bfseries СИСТЕМ ХӨГЖҮҮЛЭХ НЬ}

\vspace{1.5cm}

\begin{flushleft}
{\normalsize Хөтөлбөрийн индекс:} {\normalsize D061203}\\[4pt]
{\normalsize Хөтөлбөр:} {\normalsize Мэдээллийн системийн хөтөлбөр}\\[24pt]

{\normalsize Удирдагч:~~Дэд профессор Д.Сарангэрэл\hfill/\underline{\hspace{4cm}}/}\\[10pt]
{\normalsize Зөвлөгч:~~~Магистр Ө.Сүх-Очир\hfill/\underline{\hspace{4cm}}/}\\[10pt]
{\normalsize Гүйцэтгэгч:~Т.Оюунбаатар\hfill/\underline{\hspace{4cm}}/}
\end{flushleft}

\vfill
\begin{center}
{\large Улаанбаатар хот}\\[4pt]
{\large 2025 он}
\end{center}
\end{titlepage}
\clearpage

% ================================================================
%  ТӨЛӨВЛӨГӨӨНИЙ ХУУДАС
% ================================================================
\thispagestyle{plain}
\begin{flushleft}
{\normalsize\bfseries Батлав. Удирдагч:}\\[6pt]
\underline{\hspace{8cm}} / Дэд профессор Д.Сарангэрэл /
\end{flushleft}

\vspace{1.5cm}
\begin{center}
{\large\bfseries ТӨГСӨЛТИЙН АЖЛЫН ТӨЛӨВЛӨГӨӨ}\\[12pt]
{\normalsize Сэдэв: QR КОДООР СУУРИЛСАН РЕСТОРАНЫ ШИРЭЭНИЙ МОБАЙЛ ЗАХИАЛГЫН СИСТЕМ ХӨГЖҮҮЛЭХ НЬ}
\end{center}

\vspace{1cm}
\begin{center}
\begin{tabular}{|c|p{8cm}|c|c|}
\hline
\rowcolor{TableHdr}\textcolor{white}{\textbf{No}} &
\textcolor{white}{\textbf{Ажлын бүлэг, хэсгийн нэр}} &
\textcolor{white}{\textbf{Эзлэх хувь}} &
\textcolor{white}{\textbf{Дуусах хугацаа}} \\ \hline
1 & Удиртгал                & 5\%  & 2025-02-16 \\ \hline
2 & Онолын хэсэг            & 25\% & 2025-03-19 \\ \hline
3 & Судалгааны хэсэг        & 30\% & 2025-04-06 \\ \hline
4 & Төслийн хэсэг           & 30\% & 2025-05-04 \\ \hline
5 & Дүгнэлт                 & 10\% & 2025-05-13 \\ \hline
\end{tabular}
\end{center}

\vspace{1.5cm}
{\normalsize Төлөвлөгөөг боловсруулсан оюутан:\hfill\underline{\hspace{5cm}}/Т.Оюунбаатар/}
\clearpage

% ================================================================
%  ЗОХИОГЧ ЭРХИЙН ХАМГААЛАЛ
% ================================================================
\thispagestyle{plain}
\begin{center}
{\large\bfseries Зохиогч эрхийн хамгаалал}
\addcontentsline{toc}{chapter}{Зохиогч эрхийн хамгаалал}
\end{center}

\vspace{0.8cm}
{\normalsize
Миний бие Тогтохын Оюунбаатар, ``QR КОДООР СУУРИЛСАН РЕСТОРАНЫ ШИРЭЭНИЙ МОБАЙЛ
ЗАХИАЛГЫН СИСТЕМ ХӨГЖҮҮЛЭХ НЬ'' сэдэвт энэ ажил нь минийх бөгөөд дараахыг нотолж байна.
Үүнд:

\begin{itemize}[leftmargin=2cm]
  \item Горилогч энэ ажлыг тус сургуулиас боловсролын зэрэг авахаар бүхэлд нь буюу
        голлон хийсэн болно.
  \item Энэ ажлын аль нэг хэсгийг тус сургуульд эсвэл өөр байгууллага боловсролын
        зэрэг, мэргэшил авахаар өмнө нь илгээсэн бол түүнийгээ тодорхой заасан болно.
  \item Бусад хүмүүсийн хэвлүүлсэн ажлаас зөвлөгөө авсан бол түүнийгээ үндэслэсэн болно.
  \item Бусад хүмүүсийн ажлаас ишлэл хийхдээ гол эх үүсвэрийг нь заасан болно.
  \item Миний ажилд тусалсан голлох бүх эх үүсвэрт талархаж байна.
  \item Ажлыг бусадтай хамтарсан бол алийг нь бусад хүмүүс хийсэн болохыг тодорхой
        заасан болно.
\end{itemize}

\vspace{2cm}
Гарын үсэг: \underline{\hspace{5cm}}\\[1cm]
Огноо: \underline{\hspace{5cm}}
}
\clearpage

% ================================================================
%  ХУРААНГУЙ
% ================================================================
\thispagestyle{plain}
\begin{center}
{\large\bfseries Шинжлэх Ухаан Технологийн Их Сургууль}\\
{\large Мэдээлэл, Холбооны Технологийн Сургууль}

\vspace{0.5cm}
{\large\bfseries Хураангуй}
\addcontentsline{toc}{chapter}{Хураангуй}

\vspace{0.3cm}
{\normalsize QR кодоор суурилсан рестораны ширээний мобайл захиалгын систем
хөгжүүлэх нь}\\[6pt]
{\normalsize Т.~Оюунбаатар}\\
{\small \texttt{oyunbaatar@gmail.com}}
\end{center}

\vspace{0.8cm}
{\normalsize
QR меню ашиглан хоол захиалах систем нь ресторан, кафе болон хоолны газруудийн захиалга
өгөх процессыг бүрэн дижитал болгох зорилготой. Уламжлалт аргаар буюу зөөгчийн
тусламжтайгаар захиалга авахад зарим захиалгууд орхигдох, андуурагдах, цаг хугацаа их
зарцуулагдах хүндрэл гарч байдаг. Энэхүү ажлын хүрээнд ширээ тус бүрт байрлуулсан QR
кодыг үйлчлүүлэгч ухаалаг утсаараа уншуулж цэстэй танилцан захиалга өгөхөд тухайн
захиалга нь WebSocket протоколоор ажилтны дэлгэцэнд бодит цаг хугацаанд дамждаг
систем хэрэгжүүлсэн.

Системийг React 18 + TypeScript (вэб клиент), Node.js/Express.js (API сервер), PostgreSQL 16
(өгөгдлийн сан), Socket.IO (бодит цагийн холболт) технологиудыг ашиглан хөгжүүлсэн.
JWT токен дээр суурилсан хандалтын хяналт, Drizzle ORM-д суурилсан өгөгдлийн давхарга,
pnpm монорепо бүтцийг нэвтрүүлсэн. Туршилтын үр дүнд API хариу 120 мс дундаж, WebSocket
мессежийн дамжуулалт 45 мс, SUS (System Usability Scale) үнэлгээ 84.2 оноо, хэрэглэгч
захиалга өгөхөд зарцуулах хугацаа уламжлалт аргаас 62\% буурсан үр дүн гарсан.
}

\vspace{0.5cm}
{\small\textbf{Түлхүүр үгс:} QR код, мобайл захиалга, рестораны систем, WebSocket,
Socket.IO, бодит цагийн мэдэгдэл, React, Node.js, PostgreSQL}
\clearpage

% ================================================================
%  ТАЛАРХАЛ
% ================================================================
\thispagestyle{plain}
\begin{center}
{\large\bfseries Талархал}
\addcontentsline{toc}{chapter}{Талархал}
\end{center}

\vspace{0.8cm}
{\normalsize
Энэхүү дипломын ажлын санааг олох, хэрэгжүүлэх үр дүнг боловсруулах, нэмэлт мэдээлэл
өгөн санал солилцох зэргээр хамтран ажиллаж туслалцаа үзүүлсэн удирдагч багш
Д.~Сарангэрэл, зөвлөх багш Ө.~Сүх-Очир нартаа гүнээ талархсанаа илэрхийлье.

Мөн туршилтын явцад оролцож санал хүсэлтээ өгсөн рестораны 28 хэрэглэгч болон 4
ажилтанд талархал илэрхийлье. Судалгааны хугацаанд хамтран ажиллаж дэмжлэг үзүүлсэн
Мэдээлэл, Холбооны Технологийн Сургуулийн Мэдээллийн технологийн салбарын 2022-р
оны элсэгчдийнхээ бүлгийн гишүүдэд баярлалаа.
}
\clearpage
